% GENERATED FILE. Edit canonical lessons, then run npm run generate:books.
% canonical-source-hash: fnv1a64:4eb13986a9fc1ee7
% canonical-lessons: PA-C349-gaik, PA-C349-vanjhli, PA-C349-vaja, PA-C349-kavita, PA-C349-sarahuna

\chapter{Singers and Songs}
\label{ch:pa-a2-singers-and-songs}
\hlchaptermodality{349}

\begin{chapteropening}
\textbf{By the end of this chapter:} \emph{I can say which singers and songs I like.}

Complete the last lesson of chapter 349: I can say which singers and songs I like.
\end{chapteropening}

\section[\texorpdfstring{\pa{ਗਾਇਕ} (gāik)}{ਗਾਇਕ (gāik)}]{\pa{ਗਾਇਕ} (gāik) — a singer}

\label{lesson:PA-C349-gaik}



\begin{quote}
Before the new one: say the Punjabi for music practice, then the Punjabi for to echo.
\end{quote}

\subsection*{You'll want to know: \pa{ਗਾਇਕ}}
\textbf{\pa{ਗਾਇਕ}} — \emph{gāik} — \textquotedblleft{}a singer\textquotedblright{}.

\begin{quote}
\pa{ਉਹ} \pa{ਗਾਇਕ} \pa{ਬਹੁਤ} \pa{ਮਿੱਠਾ} \pa{ਗਾਉਂਦਾ} \pa{ਹੈ।}
\end{quote}

\begin{quote}
\emph{oh gāik bahut miṭṭhā gāundā hai.}
\end{quote}

\begin{quote}
That singer sings very sweetly.
\end{quote}

A masculine noun.

\begin{cousinweb}
\textbf{\pa{ਗਾਇਕ}} (\emph{gāik}) is taken over from Sanskrit \emph{gāyaka}, \textquotedblleft{}a singer\textquotedblright{}, a tatsama (Monier-Williams).
\end{cousinweb}

\begin{grammarlens}[title={What you've built}]
One more word for this chapter.
\end{grammarlens}

\subsection*{Your turn}
\begin{itemize}
\raggedright
  \item \emph{Say it:} \emph{gāik}
  \item \emph{Say it:} \emph{gāik}, once more
  \item \emph{Say it:} \emph{oh gāik bahut miṭṭhā gāundā hai.} — That singer sings very sweetly.
  \item \emph{Recall:} say the Punjabi for music practice, then the Punjabi for to echo, then say \emph{gāik} again
\end{itemize}

\begin{tcolorbox}[breakable,colback=teal!4,colframe=teal!35!black,title={Before you move on}]
What does \pa{ਗਾਇਕ} mean? (\textquotedblleft{}A singer\textquotedblright{}.) Say it once more.
\end{tcolorbox}
\section[\texorpdfstring{\pa{ਵੰਝਲੀ} (vanjhlī)}{ਵੰਝਲੀ (vanjhlī)}]{\pa{ਵੰਝਲੀ} (vanjhlī) — a flute; a simple pipe}

\label{lesson:PA-C349-vanjhli}



\begin{quote}
Before the new one: say the Punjabi for to echo, then the Punjabi for a singer.
\end{quote}

\subsection*{You'll want to know: \pa{ਵੰਝਲੀ}}
\textbf{\pa{ਵੰਝਲੀ}} — \emph{vanjhlī} — \textquotedblleft{}a flute; a simple pipe\textquotedblright{}.

\begin{quote}
\pa{ਮੈਨੂੰ} \pa{ਵੰਝਲੀ} \pa{ਦੀ} \pa{ਧੁਨ} \pa{ਪਸੰਦ} \pa{ਹੈ।}
\end{quote}

\begin{quote}
\emph{mainū̃ vanjhlī dī dhun pasand hai.}
\end{quote}

\begin{quote}
I like the tune of the flute.
\end{quote}

A feminine noun.

\begin{cousinweb}
\textbf{\pa{ਵੰਝਲੀ}} (\emph{vanjhlī}) is inherited from Sanskrit \emph{vaṁśī}, \textquotedblleft{}flute, pipe\textquotedblright{}; CDIAL lists Punjabi \emph{vañjhlī}, \textquotedblleft{}pipe, flute\textquotedblright{}, under it (CDIAL 11180).
\end{cousinweb}

\begin{grammarlens}[title={What you've built}]
One more word for this chapter.
\end{grammarlens}

\subsection*{Your turn}
\begin{itemize}
\raggedright
  \item \emph{Say it:} \emph{vanjhlī}
  \item \emph{Say it:} \emph{vanjhlī}, once more
  \item \emph{Say it:} \emph{mainū̃ vanjhlī dī dhun pasand hai.} — I like the tune of the flute.
  \item \emph{Recall:} say the Punjabi for to echo, then the Punjabi for a singer, then say \emph{vanjhlī} again
  \item \emph{Write it:} copy \textbf{\pa{ਵੰਝਲੀ}} once, checking the whole word
\end{itemize}

\begin{tcolorbox}[breakable,colback=teal!4,colframe=teal!35!black,title={Before you move on}]
What does \pa{ਵੰਝਲੀ} mean? (\textquotedblleft{}A flute; a simple pipe\textquotedblright{}.) Say it once more.
\end{tcolorbox}
\section[\texorpdfstring{\pa{ਵਾਜਾ} (vājā)}{ਵਾਜਾ (vājā)}]{\pa{ਵਾਜਾ} (vājā) — a harmonium; a musical instrument}

\label{lesson:PA-C349-vaja}



\begin{quote}
Before the new one: say the Punjabi for a singer, then the Punjabi for a flute.
\end{quote}

\subsection*{You'll want to know: \pa{ਵਾਜਾ}}
\textbf{\pa{ਵਾਜਾ}} — \emph{vājā} — \textquotedblleft{}a harmonium; a musical instrument\textquotedblright{}.

\begin{quote}
\pa{ਮੇਰਾ} \pa{ਦਾਦਾ} \pa{ਵਾਜਾ} \pa{ਵਜਾਉਂਦਾ} \pa{ਹੈ} \pa{ਅਤੇ} \pa{ਗਾਉਂਦਾ} \pa{ਹੈ।}
\end{quote}

\begin{quote}
\emph{merā dādā vājā vajāundā hai ate gāundā hai.}
\end{quote}

\begin{quote}
My grandfather plays the harmonium and sings.
\end{quote}

A masculine noun.

\begin{cousinweb}
\textbf{\pa{ਵਾਜਾ}} (\emph{vājā}) is inherited from Sanskrit \emph{vādya}, \textquotedblleft{}musical instrument\textquotedblright{}; CDIAL lists Punjabi \emph{vājjā}, \textquotedblleft{}musical instrument\textquotedblright{}, under it (CDIAL 11511). CDIAL does not give the narrower sense \textquotedblleft{}harmonium\textquotedblright{}.
\end{cousinweb}

\begin{grammarlens}[title={What you've built}]
One more word for this chapter.
\end{grammarlens}

\subsection*{Your turn}
\begin{itemize}
\raggedright
  \item \emph{Say it:} \emph{vājā}
  \item \emph{Say it:} \emph{vājā}, once more
  \item \emph{Say it:} \emph{merā dādā vājā vajāundā hai ate gāundā hai.} — My grandfather plays the harmonium and sings.
  \item \emph{Recall:} say the Punjabi for a singer, then the Punjabi for a flute, then say \emph{vājā} again
  \item \emph{Write it:} copy \textbf{\pa{ਵਾਜਾ}} once, checking the whole word
\end{itemize}

\begin{tcolorbox}[breakable,colback=teal!4,colframe=teal!35!black,title={Before you move on}]
What does \pa{ਵਾਜਾ} mean? (\textquotedblleft{}A harmonium; a musical instrument\textquotedblright{}.) Say it once more.
\end{tcolorbox}
\section[\texorpdfstring{\pa{ਕਵਿਤਾ} (kavitā)}{ਕਵਿਤਾ (kavitā)}]{\pa{ਕਵਿਤਾ} (kavitā) — a poem; poetry}

\label{lesson:PA-C349-kavita}



\begin{quote}
Before the new one: say the Punjabi for a flute, then the Punjabi for a harmonium.
\end{quote}

\subsection*{You'll want to know: \pa{ਕਵਿਤਾ}}
\textbf{\pa{ਕਵਿਤਾ}} — \emph{kavitā} — \textquotedblleft{}a poem; poetry\textquotedblright{}.

\begin{quote}
\pa{ਮੇਰੀ} \pa{ਭੈਣ} \pa{ਹਰ} \pa{ਹਫ਼ਤੇ} \pa{ਇੱਕ} \pa{ਕਵਿਤਾ} \pa{ਲਿਖਦੀ} \pa{ਹੈ।}
\end{quote}

\begin{quote}
\emph{merī bhaiṇ har hafte ikk kavitā likhdī hai.}
\end{quote}

\begin{quote}
My sister writes a poem every week.
\end{quote}

A feminine noun.

\begin{cousinweb}
\textbf{\pa{ਕਵਿਤਾ}} (\emph{kavitā}) is taken over from Sanskrit \emph{kavitā}, \textquotedblleft{}a poem\textquotedblright{}, a tatsama (Monier-Williams).
\end{cousinweb}

\begin{grammarlens}[title={What you've built}]
One more word for this chapter.
\end{grammarlens}

\subsection*{Your turn}
\begin{itemize}
\raggedright
  \item \emph{Say it:} \emph{kavitā}
  \item \emph{Say it:} \emph{kavitā}, once more
  \item \emph{Say it:} \emph{merī bhaiṇ har hafte ikk kavitā likhdī hai.} — My sister writes a poem every week.
  \item \emph{Recall:} say the Punjabi for a flute, then the Punjabi for a harmonium, then say \emph{kavitā} again
  \item \emph{Write it:} copy \textbf{\pa{ਕਵਿਤਾ}} once, checking the whole word
\end{itemize}

\begin{tcolorbox}[breakable,colback=teal!4,colframe=teal!35!black,title={Before you move on}]
What does \pa{ਕਵਿਤਾ} mean? (\textquotedblleft{}A poem; poetry\textquotedblright{}.) Say it once more.
\end{tcolorbox}
\section[\texorpdfstring{\pa{ਸਰਾਹੁਣਾ} (sarāhuṇā)}{ਸਰਾਹੁਣਾ (sarāhuṇā)}]{\pa{ਸਰਾਹੁਣਾ} (sarāhuṇā) — to praise}

\label{lesson:PA-C349-sarahuna}



\begin{quote}
Before the new one: say the Punjabi for a harmonium, then the Punjabi for a poem.
\end{quote}

\subsection*{You'll want to know: \pa{ਸਰਾਹੁਣਾ}}
\textbf{\pa{ਸਰਾਹੁਣਾ}} — \emph{sarāhuṇā} — \textquotedblleft{}to praise\textquotedblright{}.

\begin{quote}
\pa{ਲੋਕ} \pa{ਉਸ} \pa{ਗਾਇਕ} \pa{ਨੂੰ} \pa{ਬਹੁਤ} \pa{ਸਰਾਹੁੰਦੇ} \pa{ਹਨ।}
\end{quote}

\begin{quote}
\emph{lok us gāik nū̃ bahut sarāhunde han.}
\end{quote}

\begin{quote}
People praise that singer a great deal.
\end{quote}

Past (I …): \pa{ਮੈਂ} \pa{ਉਸ} \pa{ਦੀ} \pa{ਕਵਿਤਾ} \pa{ਸਰਾਹੀ}, \emph{maiṁ us dī kavitā sarāhī}.

\begin{cousinweb}
\textbf{\pa{ਸਰਾਹੁਣਾ}} (\emph{sarāhuṇā}) is inherited from Sanskrit \emph{ślāghate}, \textquotedblleft{}praises\textquotedblright{}; CDIAL lists Punjabi \emph{sarāhuṇā}, \textquotedblleft{}praises\textquotedblright{}, under it (CDIAL 12734).
\end{cousinweb}

\begin{grammarlens}[title={What you've built}]
One more word for this chapter.
\end{grammarlens}

\subsection*{Your turn}
\begin{itemize}
\raggedright
  \item \emph{Say it:} \emph{sarāhuṇā}
  \item \emph{Say it:} \emph{sarāhuṇā}, once more
  \item \emph{Say it:} \emph{lok us gāik nū̃ bahut sarāhunde han.} — People praise that singer a great deal.
  \item \emph{Recall:} say the Punjabi for a harmonium, then the Punjabi for a poem, then say \emph{sarāhuṇā} again
  \item \emph{Write it:} copy \textbf{\pa{ਸਰਾਹੁਣਾ}} once, checking the whole word
\end{itemize}

\begin{tcolorbox}[breakable,colback=teal!4,colframe=teal!35!black,title={Before you move on}]
What does \pa{ਸਰਾਹੁਣਾ} mean? (\textquotedblleft{}To praise\textquotedblright{}.) Say it once more.
\end{tcolorbox}
